\section*{Capítulo \ref{ch:b1:respiration-fermentation} --- Respiración celular y fermentación}

\begin{solution}{exo:b1:respiration-fermentation:1}
\omterm{def:b1:respiration-fermentation:glycolysis}{Glucólisis}, \omterm{def:b1:eukaryotic-cell:organelle}{citosol}: 2 piruvatos, 2 \omterm{def:b1:water-small-molecules:atp}{ATP}, 2 NADH. Oxidación del piruvato y
\omterm{def:b1:respiration-fermentation:krebs}{ciclo de Krebs}, \omterm{def:b1:eukaryotic-cell:mitochondrion}{matriz mitocondrial}: 6 $\mathrm{CO_2}$, 8 NADH, 2
$\mathrm{FADH_2}$, 2 GTP. Fosforilación oxidativa, membrana interna: agua a
partir del oxígeno, unos 26 \omterm{def:b1:water-small-molecules:atp}{ATP}.
\end{solution}

\begin{solution}{exo:b1:respiration-fermentation:2}
Glucosa $+ 2\,\mathrm{NAD^+} + 2\,\mathrm{ADP} + 2\,\mathrm{P_i} \to
2$ piruvatos $+ 2\,\mathrm{NADH} + 2\,\mathrm{H^+} + 2\,\mathrm{ATP} +
2\,\mathrm{H_2O}$. \omterm{def:b1:respiration-fermentation:glycolysis}{Fosforilación a nivel de sustrato}: una \omterm{def:b1:enzymes:enzyme}{enzima} transfiere
un fosfato directamente desde un intermediario rico en energía
(1,3-bisfosfoglicerato, fosfoenolpiruvato) al ADP, sin gradiente de membrana
ni oxígeno.
\end{solution}

\begin{solution}{exo:b1:respiration-fermentation:3}
NADH y $\mathrm{CO_2}$ (isocitrato a $\alpha$-cetoglutarato), NADH y
$\mathrm{CO_2}$ ($\alpha$-cetoglutarato a succinil-CoA), GTP (succinil-CoA a
succinato), $\mathrm{FADH_2}$ (succinato a fumarato), NADH (malato a
oxalacetato). Los dos $\mathrm{CO_2}$ proceden de las dos descarboxilaciones
de los ácidos de seis y cinco carbonos.
\end{solution}

\begin{solution}{exo:b1:respiration-fermentation:4}
Regenera el $\mathrm{NAD^+}$ para que la \omterm{def:b1:respiration-fermentation:glycolysis}{glucólisis} pueda continuar y rendir
sus 2 \omterm{def:b1:water-small-molecules:atp}{ATP} por glucosa. No oxida el carbono, ni extrae la energía restante
(el producto la conserva casi entera), ni fabrica más de la quinceava parte
del \omterm{def:b1:water-small-molecules:atp}{ATP} de la respiración.
\end{solution}

\begin{solution}{exo:b1:respiration-fermentation:5}
Al oxígeno: $-2\times 96\,500\times 1.14 = \qty{-220}{kJ/mol}$, cuatro \omterm{def:b1:water-small-molecules:atp}{ATP} a
\qty{50}{kJ}. A la ubiquinona: $-2\times 96\,500\times 0.37 =
\qty{-71}{kJ/mol}$, un \omterm{def:b1:water-small-molecules:atp}{ATP}.
\end{solution}

\begin{solution}{exo:b1:respiration-fermentation:6}
Músculo: $2 + 2$ (GTP) $+ 8\times 2.5$ (NADH mitocondrial) $+ 2\times
1.5$ ($\mathrm{FADH_2}$) $+ 2\times 1.5$ (NADH citosólico vía \omterm{def:b1:respiration-fermentation:carriers}{FAD}) $=
30$. Hígado: el NADH citosólico da $2\times 2.5$: 32.
\end{solution}

\begin{solution}{exo:b1:respiration-fermentation:7}
CR $= 2.2$. Glucosa respirada $x$: $6x = 1.0$, $x = \qty{0.167}{mmol}$;
$\mathrm{CO_2}$: $6x + 2y = 2.2$, $2y = 1.2$, $y = \qty{0.6}{mmol}$. Se
fermenta $0.6/0.767 = \qty{78}{\%}$ de la glucosa y se respira el
\qty{22}{\%}.
\end{solution}

\begin{solution}{exo:b1:respiration-fermentation:8}
El consumo de oxígeno sube (la cadena, liberada de la contrapresión del
gradiente, funciona a toda velocidad), la síntesis de \omterm{def:b1:water-small-molecules:atp}{ATP} se detiene (no hay
gradiente que impulse la sintasa) y toda la energía de los electrones
aparece como calor. Los pacientes quemaban su \omterm{def:b1:lipids:triglyceride}{grasa} deprisa y morían de una
temperatura corporal que nada podía bajar.
\end{solution}

\begin{solution}{exo:b1:respiration-fermentation:9}
La cadena se atasca: todos los transportadores quedan reducidos, no se
bombean protones y no se fabrica \omterm{def:b1:water-small-molecules:atp}{ATP}. El NADH no se vuelve a oxidar, de modo
que el \omterm{def:b1:respiration-fermentation:krebs}{ciclo de Krebs} y la piruvato deshidrogenasa se detienen por falta de
$\mathrm{NAD^+}$. La \omterm{def:b1:respiration-fermentation:glycolysis}{glucólisis}, liberada de la inhibición por \omterm{def:b1:water-small-molecules:atp}{ATP}, funciona
deprisa y debe regenerar su $\mathrm{NAD^+}$ reduciendo el piruvato a
lactato, que inunda la sangre.
\end{solution}

\begin{solution}{exo:b1:respiration-fermentation:10}
$8\times 10 + 7\times 2.5 + 7\times 1.5 - 2 = 80 + 17.5 + 10.5 - 2 =
106$ \omterm{def:b1:water-small-molecules:atp}{ATP}; $106/16 = 6.6$ por carbono. Oxígeno: el palmitato, 23
$\mathrm{O_2}$ para 106 \omterm{def:b1:water-small-molecules:atp}{ATP}, \qty{0.22}{O_2} por \omterm{def:b1:water-small-molecules:atp}{ATP}; la glucosa, 6 para 30,
\qty{0.20}{O_2} por \omterm{def:b1:water-small-molecules:atp}{ATP}. Por unidad de oxígeno la glucosa es un
\qty{10}{\%} mejor: la foca, limitada por el oxígeno bajo el agua, prefiere
los \omterm{def:b1:carbohydrates:monosaccharide}{glúcidos}; por unidad de masa la \omterm{def:b1:lipids:triglyceride}{grasa} es el doble de buena: el colibrí,
limitado por la masa en vuelo, migra con \omterm{def:b1:lipids:triglyceride}{grasa}.
\end{solution}

\begin{solution}{exo:b1:respiration-fermentation:11}
Solo el \omterm{def:b1:water-small-molecules:atp}{ATP}: $5/3 < \qty{2}{s}$. Fosfocreatina: $25/3 = \qty{8}{s}$.
\omterm{def:b1:respiration-fermentation:fermentation}{Fermentación}: $2\times 3 = \qty{6}{mmol/L}$ de \omterm{def:b1:water-small-molecules:atp}{ATP} por segundo, suficiente
para la demanda, mientras duren el \omterm{def:b1:carbohydrates:polysaccharide}{glucógeno} y la tolerancia al lactato
(decenas de segundos). El aporte de oxígeno tarda un minuto o más en subir y
suministra como mucho \qty{1}{mmol/L} de \omterm{def:b1:water-small-molecules:atp}{ATP} por segundo en una fibra: una
carrera de diez segundos termina antes de que la respiración haya empezado a
ayudar, y la pagan la fosfocreatina y la \omterm{def:b1:respiration-fermentation:fermentation}{fermentación}.
\end{solution}

\begin{solution}{exo:b1:respiration-fermentation:12}
El gradiente almacena energía como una diferencia de concentración de
protones y de potencial eléctrico a través de una membrana aislante ---
exactamente un condensador cargado más una pila de concentración ---; la
cadena es el cargador, impulsado por la caída de los electrones; y la
sintasa es un motor que hace girar la corriente de protones, exactamente una
turbina, que acopla el flujo a una rotación mecánica que fabrica \omterm{def:b1:water-small-molecules:atp}{ATP}. La
analogía es vaga en que la ``batería'' contiene solo milisegundos del
consumo de la \omterm{def:b1:cell-unit-of-life:cell}{célula} y debe cargarse sin cesar, en que los demás
\omterm{def:b1:membranes-transport:transporters}{transportadores} de esa misma membrana también recurren a ella, y en que la
turbina puede girar al revés e hidrolizar \omterm{def:b1:water-small-molecules:atp}{ATP} para bombear protones cuando
la cadena falla.
\end{solution}

\begin{solution}{pb:b1:respiration-fermentation:1}
\textbf{1.} \omterm{def:b1:respiration-fermentation:glycolysis}{Glucólisis}: 2 \omterm{def:b1:water-small-molecules:atp}{ATP}, 2 NADH (citosólicos). Piruvato
deshidrogenasa: 2 NADH. Krebs: 6 NADH, 2 $\mathrm{FADH_2}$, 2 GTP.
\textbf{2.} NADH mitocondrial, $8\times 10 = 80$; $\mathrm{FADH_2}$,
$2\times 6 = 12$; NADH citosólico como $\mathrm{FADH_2}$, $2\times 6 = 12$:
104 protones.
\textbf{3.} $104/4 = 26$ \omterm{def:b1:water-small-molecules:atp}{ATP}; total, $26 + 2 + 2 = 30$.
\textbf{4.} 2 \omterm{def:b1:water-small-molecules:atp}{ATP}.
\textbf{5.} 15.
\textbf{6.} Con aire, $30\times 50 = \qty{1500}{kJ}$: el \qty{52}{\%} de
2870. Sin aire: \qty{100}{kJ}, el \qty{3.5}{\%} de 2870, pero el
\qty{43}{\%} de los 235 liberados; la \omterm{def:b1:respiration-fermentation:fermentation}{fermentación} es eficaz capturando lo
que libera, solo que libera poco.
\textbf{7.} Aireado, $30/30 = \qty{1}{mmol}$ de glucosa por hora; cerrado,
$30/2 = \qty{15}{mmol}$.
\textbf{8.} Cerrado: \qty{30}{mmol} de etanol y \qty{30}{mmol} de
$\mathrm{CO_2}$ por hora. Aireado: \qty{6}{mmol} de $\mathrm{CO_2}$
liberados y \qty{6}{mmol} de $\mathrm{O_2}$ consumidos.
\textbf{9.} $2\times 10.5 = \qty{21}{g}$ por mol, \qty{0.12}{g} por gramo de
glucosa.
\textbf{10.} $0.5\times 180 = \qty{90}{g}$ por mol, es decir, $90/10.5 =
8.6$ \omterm{def:b1:water-small-molecules:atp}{ATP} de equivalente frente a los 30 disponibles: el \omterm{def:b1:water-small-molecules:atp}{ATP} sobra; el
rendimiento lo limita el carbono (la mitad de la glucosa se oxida a
$\mathrm{CO_2}$ para fabricar el \omterm{def:b1:water-small-molecules:atp}{ATP}, y la biomasa necesita esqueletos
carbonados y poder reductor), y el \omterm{def:b1:water-small-molecules:atp}{ATP} sobrante se gasta en mantenimiento.
\textbf{11.} Para obtener el mismo \omterm{def:b1:water-small-molecules:atp}{ATP}, el cultivo cerrado consume quince
veces más glucosa (pregunta 7) y crece una quinta parte por gramo (pregunta
9): consumo rápido de azúcar, poco crecimiento y alcohol; la observación de
Pasteur. La fosfofructoquinasa, inhibida por el \omterm{def:b1:water-small-molecules:atp}{ATP} y el citrato con
aireación, queda liberada cuando cesa la respiración.
\textbf{12.} $\mathrm{CO_2}$, $44/44 = \qty{1.0}{mmol}$; $\mathrm{O_2}$,
$20/32 = \qty{0.625}{mmol}$: CR $= 1.6$.
\textbf{13.} $\mathrm{O_2}$: $6x = 0.625$, $x = \qty{0.104}{mmol}$.
$\mathrm{CO_2}$: $6x + 2y = 1.0$, $2y = 0.375$, $y = \qty{0.1875}{mmol}$.
\textbf{14.} Se fermenta $0.1875/0.292 = \qty{64}{\%}$ de la glucosa. El \omterm{def:b1:water-small-molecules:atp}{ATP}
de la \omterm{def:b1:respiration-fermentation:fermentation}{fermentación} es $2y = 0.375$ frente a $30x = 3.12$ de la respiración:
el \qty{11}{\%}.
\textbf{15.} $16/23 = 0.70$.
\textbf{16.} En ayunas, el animal quema \omterm{def:b1:lipids:triglyceride}{grasa} (CR 0,7); alimentado, quema
los \omterm{def:b1:carbohydrates:monosaccharide}{glúcidos} de la comida (1,0).
\textbf{17.} $96\,500\times 0.22 = \qty{21.2}{kJ/mol}$.
\textbf{18.} $4\times 21.2 = \qty{85}{kJ}$ por cada \omterm{def:b1:water-small-molecules:atp}{ATP} de \qty{50}{kJ}:
un \qty{59}{\%}.
\textbf{19.} $2\times 96\,500\times 1.14 = \qty{220}{kJ}$; diez protones
almacenan $10\times 21.2 = \qty{212}{kJ}$: un \qty{96}{\%}; la cadena
desperdicia poco, y las pérdidas están en la sintasa y en los
\omterm{def:b1:membranes-transport:transporters}{transportadores}.
\textbf{20.} $0.96\times 0.59 = 0.57$; $2.5\times 50/220 = 0.57$: la misma
cifra por los dos caminos.
\textbf{21.} El consumo de oxígeno sube, el rendimiento en \omterm{def:b1:water-small-molecules:atp}{ATP} cae hacia el
de la \omterm{def:b1:respiration-fermentation:fermentation}{fermentación} (2 por glucosa de la \omterm{def:b1:respiration-fermentation:glycolysis}{glucólisis}, y como la PFK queda
liberada, el consumo de glucosa sube), la producción de calor sube y el
crecimiento se desploma.
\textbf{22.} La síntesis de \omterm{def:b1:water-small-molecules:atp}{ATP} se detiene; el gradiente no puede
descargarse, la cadena se atasca contra él, el consumo de oxígeno cesa y la
\omterm{def:b1:cell-unit-of-life:cell}{célula} fermenta lo que puede. Añadir después un desacoplante deja que los
protones vuelvan a entrar: la cadena y el consumo de oxígeno se reanudan,
todavía sin \omterm{def:b1:water-small-molecules:atp}{ATP}.
\textbf{23.} El NADH solo bombea por los complejos III y IV: $6/4 = 1.5$;
por glucosa, $8\times 1.5 + 2\times 1.5 + 2\times 1.5 + 4 = 22$ \omterm{def:b1:water-small-molecules:atp}{ATP}.
\textbf{24.} Sin oxígeno la cadena no puede aceptar electrones, así que el
NADH no se vuelve a oxidar y la reserva permanece reducida; las tres
deshidrogenasas dependientes de \omterm{def:b1:respiration-fermentation:carriers}{NAD} del ciclo se quedan sin
$\mathrm{NAD^+}$ y el ciclo se detiene, que es por lo que la \omterm{def:b1:respiration-fermentation:fermentation}{fermentación}
debe regenerar $\mathrm{NAD^+}$ en el \omterm{def:b1:eukaryotic-cell:organelle}{citosol} para que la \omterm{def:b1:respiration-fermentation:glycolysis}{glucólisis} siga.
\textbf{25.} 30 \omterm{def:b1:water-small-molecules:atp}{ATP} con aire y 2 sin él: un cociente de 15; P/O de 2,5 para
el NADH (10 protones sobre 4 por \omterm{def:b1:water-small-molecules:atp}{ATP}) y de 1,5 para el $\mathrm{FADH_2}$ (6
sobre 4).
\end{solution}
